% =====================================================================
%  DRIFTGUARD — Section IV: Experimental Evaluation.
%  Narrative rewrite (NOT a clean-up of Experimental_Evaluation.tex).
%  Numbers are taken verbatim from the bible; prose is fresh, story-first.
%  Six tables kept in-body: data, leakage, xdataset, oracle, al-strategy,
%  cursor. The alpha-sensitivity table are deferred to
%  the appendix; only their takeaways appear here.
% =====================================================================

\section{Experimental Evaluation}
\label{sec:eval}

Our evaluation follows the life of the system and answers four questions in
order. Does the leakage gate buy us an \emph{honest} baseline, one whose
cross-distribution accuracy reflects phishing rather than collection artefacts
(Section~\ref{ssec:eval-gate})? Given that baseline, which evidence makes the
LLM the best possible oracle (Section~\ref{ssec:eval-oracle})? How much drift
does the self-repair loop actually recover, at what operational cost, and
without forgetting (Section~\ref{ssec:eval-recovery})? And finally, why is an
\emph{imperfect} oracle not merely tolerable but causally \emph{preferable} for
robustness (Section~\ref{sec:noise})? Unless noted, multi-seed results are
reported over five clean seeds as mean $\pm$ 95\% confidence interval;
single-stream operational analyses use the representative seed $2025$.

\subsection{Datasets and Protocol}
\label{ssec:eval-data}

We use three datasets, each with a fixed and distinct role
(Table~\ref{tab:data}). The detector is trained \emph{only} on KnowPhish (KP).
PhreshPhish (PP), collected in 2024--25, is a temporally disjoint corpus of
fresh phishing and is our principal drift test set. DeltaPhish/SpacePhish (DP)
is held out entirely as an adversarial out-of-distribution (OOD) probe and is
\emph{never} trained on. The two kept apart deliberately: KP$\to$PP measures
natural temporal drift, KP$\to$DP measures adversarial robustness, and conflating
them is exactly how prior work overstates generalisation.

For the self-repair loop, KP is split $80/20$ into a training fold (carried
unchanged into every retrain) and an in-distribution test (the forgetting
check). PP is replayed as a row-by-row \emph{stream} in chronological order ---
so the simulation sees drift as it happened --- with a disjoint held-out test
set the oracle never touches. Three invariants are enforced in code and matter
for every number that follows: DP labels never enter training, the PP test set
is isolated from both the stream and the labelling pool, and the drift detector
sees only the model's own output, never a ground-truth label.

\begin{table}[H]
\centering
\caption{The data and its operational slices ($\mathrm{seed}=2025$). Only
$\mathrm{KP_{train}}$ and accepted oracle labels are ever trained on; DP is a
pure OOD probe.}
\label{tab:data}
\small
\begin{tabular}{@{}llrl@{}}
\toprule
Slice & Source & $N$ & Role \\
\midrule
$\mathrm{KP_{train}}$ & KP 80\% & 7{,}221 & seed training set (frozen) \\
$\mathrm{KP_{test}}$  & KP 20\% & 1{,}806 & forgetting check \\
$\mathrm{PP_{stream}}$ & PP 80\% & 3{,}821 & drift stream (1{,}865 phish) \\
$\mathrm{PP_{test}}$  & PP 20\% & 956 & drift-recovery metric (held out) \\
$\mathrm{DP_{probe}}$ & DP full & 6{,}426 & OOD probe (85\% benign, eval only) \\
\bottomrule
\end{tabular}
\end{table}

\subsection{An Honest Baseline}
\label{ssec:eval-gate}

The leakage gate rejects $13$ of $39$ candidate features, and the rejections
tell a coherent story (Table~\ref{tab:leakage}). The archival and popularity
signals --- Wayback, Tranco, DNS --- are near-perfect on KP yet no better than a
coin flip on DP: as anticipated in Section~\ref{ssec:core}, they encode
``alive versus dead'', a property of \emph{when} the data was collected, not of
phishing. The same rule, with no special-casing, also rejects the proposers'
own headline probabilities (\texttt{max\_prob}, \texttt{prob\_sn}): brilliant on
KP, useless on fresh phishing. That a single statistical criterion catches a
shortcut in the base detectors and a temporal artefact in the OSINT layer ---
two failures of entirely different origin --- is the methodological payoff of
this stage. What survives is consistent across datasets: of the proposer
signals only \texttt{min\_prob} and \texttt{prob\_pg} generalise, joined by
registration-age features that are robust because domain age is genuinely
informative.

\begin{table}[H]
\centering
\caption{Representative gate decisions. AUC is the polarity-corrected
single-feature AUC per dataset. The same rule rejects temporally leaky OSINT
signals \emph{and} the proposers' non-transferable probabilities.}
\label{tab:leakage}
\small
\begin{tabular}{@{}lcccl@{}}
\toprule
Feature & $\mathrm{AUC}_{kp}$ & $\mathrm{AUC}_{dp}$ & $\mathrm{AUC}_{pp}$ & Verdict \\
\midrule
\texttt{wayback\_n\_snapshots}   & 0.996 & 0.504 & 0.914 & reject \\
\texttt{max\_prob}               & 0.989 & 0.802 & 0.504 & reject \\
\texttt{prob\_sn}                & 0.978 & 0.761 & 0.513 & reject \\
\texttt{tranco\_rank}            & 0.962 & 0.500 & 0.655 & reject \\
\midrule
\texttt{min\_prob}               & 0.982 & 0.978 & 0.858 & keep \\
\texttt{prob\_pg}                & 0.976 & 0.987 & 0.863 & keep \\
\texttt{rdap\_domain\_age\_days} & 0.884 & 0.503 & 0.859 & keep \\
\bottomrule
\end{tabular}
\end{table}

Sterilising the feature space costs almost nothing in-distribution --- all three
candidate estimators still exceed $98\%$ weighted $F_1$ on KP --- but changes the
picture out of distribution entirely (Table~\ref{tab:xdataset}). Ungated, the
meta-learner degenerates on the adversarial probe, calling almost everything
phishing ($F_{1,w}=7.1\%$); gated, that pathology is gone. Among the three
estimators, logistic regression maximises PP, the drifted distribution we care
about, and is the deployed detector. It begins where a good self-repair target
should: high phishing precision ($0.843$) and low recall ($0.372$). The pages it
calls phishing are almost always phishing; the ones it misses are exactly the
territory the loop must reclaim. We adopt its $\mathrm{PP}\,F_{1,w}=0.629$ as the
leakage-clean baseline. This is, by design, \emph{below} the stronger base
detector taken alone (PhishGraph reaches $0.748$ on PP) --- we have thrown away
the shortcuts that inflate it --- and the contribution is not a higher starting
point but a system that, after adaptation, climbs past the static ceiling to
$0.791$ while a frozen detector cannot move at all.

\begin{table}[H]
\centering
\caption{Weighted $F_1$ per split under the gated feature space. LogReg
maximises PP, the drift target, and is deployed. Ungated, the same pipeline
collapses to $7.1\%$ on the adversarial DP probe.}
\label{tab:xdataset}
\small
\begin{tabular}{@{}lccc@{}}
\toprule
Model & KP test & DP full & PP full \\
\midrule
LogReg & 0.981 & 0.514 & \textbf{0.629} \\
GB     & \textbf{0.990} & 0.481 & 0.592 \\
RF     & 0.989 & \textbf{0.535} & 0.592 \\
\bottomrule
\end{tabular}
\end{table}

\subsection{Seeing vs.\ Cheating: A Teacher--Student Gap}
\label{ssec:eval-gap}

The gate makes an asymmetric choice: it strips eleven lifecycle features ---
Tranco, Wayback, DNS, ASN --- from the detector, yet forwards them to the oracle
(Section~\ref{ssec:core}). Is that the right split, or does it needlessly blind
the detector? Three observations show the asymmetry is exactly right and, in the
process, reveal a genuine teacher--student gap that is a secondary message of
this paper.

First, the oracle truly leans on this evidence. Across $4{,}485$ cached verdicts,
$93\%$ of its stated reasoning invokes signals the detector cannot represent, and
two in three explicitly name Tranco or Wayback: the oracle is reasoning about a
domain's lifecycle, not only its page.

Second, those same signals are \emph{leakage} for a statistical model.
\texttt{wayback\_n\_snapshots} scores AUC $0.996$ on KP and $0.914$ on PP but
exactly chance ($0.504$) on the adversarial probe (Table~\ref{tab:leakage}): it
tracks ``alive versus dead'', which lines up with the label only because of when
each corpus was collected.

Third --- the crux --- if we hand the eleven features back to the detector, it
does not learn what the oracle knows; it learns to cheat
(Figure~\ref{fig:restoration}). The detector seizes on them eagerly: they take
$37\%$ of its weight and \texttt{wayback\_n\_snapshots} becomes its single most
important feature. But the gain is leakage, not understanding. In a controlled
replay on a fixed sample path, the restored detector reaches
$\mathrm{PP}\,F_{1,w}=0.899$ --- \emph{above} the $0.859$ the very same pipeline
attains with perfect ground-truth labels (and above the main loop's $0.873$
ceiling of Section~\ref{sec:noise}) --- a level only the temporal artefact can
buy, while KP saturates from $0.968$ to $0.996$. The lift lands precisely on the
time-correlated corpora and corresponds to nothing the model could generalise
with: on the adversarial probe the same feature is chance-level
(Table~\ref{tab:leakage}), so there is no shortcut to take.

\begin{figure}[t]
  \centering
  \includegraphics[width=\columnwidth]{figures/fig_restoration.pdf}
  \caption{The student can only cheat with the teacher's features. Handed the
  eleven gated-out OSINT signals, the statistical detector reaches a PP score
  (restored, $0.899$) \emph{beyond} what the same pipeline achieves with perfect
  ground-truth labels (the $0.859$ ceiling) --- impossible without exploiting the
  temporal leakage the gate removed --- and saturates KP. The gain is confined to
  the time-correlated corpora; on the adversarial probe the same feature is
  chance-level. Absolute values are from a controlled fixed-path replay and are
  not comparable to the main loop (Table~\ref{tab:al-strategy}); the
  gated-vs-restored contrast and the ceiling crossing are the point.}
  \label{fig:restoration}
\end{figure}

The picture is now coherent. Detector and oracle see the same signal but can use
it in qualitatively different ways. A statistical model can use it only as a
correlation, and that correlation \emph{is} the temporal leakage the gate
removed; the oracle reads it in context --- ``recently registered, never
archived, absent from Tranco'' together spell throwaway infrastructure --- and so
generalises. The teacher--student gap is therefore not that the student lacks
features: it is that the only way the student \emph{can} use them is the one the
gate rightly forbids. The student does not fail to \emph{see}; it can only
\emph{cheat} --- which is exactly why the same lifecycle evidence that would
corrupt the detector is what makes the oracle worth querying.

\subsection{The Right Evidence for the Oracle}
\label{ssec:eval-oracle}

Before wiring the oracle into the loop we ask what to show it, on the cases that
matter: the $100$ most diagnostic detector errors on PP --- the $50$ phishing it
was most sure were benign and the $50$ benign it was most sure were phishing ---
so that every sample carries a detector probability that is confidently
\emph{wrong}. The result is counter-intuitive: \textbf{handing the oracle the
detector's probabilities makes it worse} (Table~\ref{tab:oracle}). Removing them
lifts weighted $F_1$ by $22$ points and raises recovery of missed phishing from
$24/50$ to $40/50$. The mechanism is plain once stated --- on these samples the
probability is a confidently wrong prior, and the model partially anchors to it
--- and it is, read positively, a robustness property: the factual OSINT dossier
gives the oracle independent ground to stand on.

\begin{table}[H]
\centering
\caption{Oracle ablation on $100$ boundary samples ($50$ FN, $50$ FP). FN/FP
recovery is the fraction of the detector's errors the oracle corrects.
Probability-blind variants (C, D) dominate on these confidently-wrong cases.}
\label{tab:oracle}
\small
\begin{tabular}{@{}lccccc@{}}
\toprule
Variant & Acc & $F_{1,w}$ & $R_{\text{phish}}$ & FN rec. & FP rec. \\
\midrule
D \texttt{html\_osint}        & \textbf{0.910} & \textbf{0.909} & 0.820 & 41/50 & 50/50 \\
C \texttt{html\_only}         & 0.900 & 0.899 & 0.800 & 40/50 & 50/50 \\
E \texttt{poc\_2k}            & 0.820 & 0.814 & 0.640 & 32/50 & 50/50 \\
B \texttt{html\_probs\_osint} & 0.770 & 0.761 & 0.580 & 28/50 & 48/50 \\
A \texttt{html\_probs}        & 0.690 & 0.676 & 0.480 & 24/50 & 45/50 \\
\bottomrule
\end{tabular}
\end{table}

There is a twist that governs deployment. This test set is, by construction,
adversarial to the probabilities, which is why the probability-blind variants
win on it. The samples the loop actually queries are different: they sit near
the genuine decision boundary, where a probability encodes ``uncertain'' rather
than ``confidently wrong'', and there the relationship inverts and the full
multimodal variant becomes the stronger oracle. The lesson is that the best
oracle context is distribution-dependent --- probabilities help when they encode
real uncertainty and hurt when they encode confident error --- and we accordingly
adopt the full multimodal variant inside the loop, where its samples live.

\subsection{Recovery, Forgetting, and Selection}
\label{ssec:eval-recovery}

This is the main result (Table~\ref{tab:al-strategy}). Over five clean seeds the
self-repair loop lifts $\mathrm{PP_{test}}$ weighted $F_1$ from $0.623$ to
$\mathbf{0.791 \pm 0.013}$ --- about $+17$ points --- and it does so in the right
way: the gain is driven entirely by phishing recall ($0.37 \to 0.61$) while
precision is \emph{preserved} ($\approx 0.96$). The loop recovers missed
phishing without sacrificing the high-precision operating point that made the
detector trustworthy. And it does so without forgetting: across all $40$ runs the
in-distribution KP score never leaves a one-point band ($0.972$--$0.982$), the
empirical signature of full-union retraining. At roughly one to two dollars of
tokens per run, this is the headline the paper promises.

We compare our \emph{drift-anchored} selector (Section~\ref{sec:design}) against
the standard active-learning acquisition rules, each of which decides which pooled
samples to send to the oracle. \emph{Margin} is classical uncertainty sampling:
the samples whose predicted probability sits closest to the $0.5$ boundary.
\emph{QBC} (query-by-committee) picks the samples on which a committee of models
most disagrees. \emph{Core-set} ignores uncertainty and instead chooses a diverse
subset that geometrically covers the feature space (a $k$-center rule).
\emph{BADGE} fuses the two, clustering gradient embeddings that are large when the
model is uncertain and spread out when samples are diverse. \emph{Hybrid} is a
fixed blend of margin, core-set diversity, and committee disagreement.
\emph{Random} draws uniformly from the pool --- the budget-matched control that
isolates how much the \emph{choice} of samples matters beyond their number.
Finally, the \emph{ungated stream} row removes the uncertainty pool entirely and
labels the raw delegate stream, isolating the pool's contribution.

\begin{table}[H]
\centering
\caption{Selection strategies over five seeds (periodic trigger, multimodal
oracle). PP is weighted $F_1$; DP is macro $F_1$ (the probe is $85\%$ benign);
``oracle acc'' is the agreement of accepted labels with ground truth. Round-0
baselines: $\mathrm{PP}=0.623$, $\mathrm{DP}=0.438$. The informed strategies are
indistinguishable on PP but separate sharply on the OOD probe.}
\label{tab:al-strategy}
\small
\begin{tabular}{@{}lccc@{}}
\toprule
Strategy & $\mathrm{PP}\,F_{1,w}$ & $\mathrm{DP}\,F_{1,\mathrm{mac}}$ & oracle acc \\
\midrule
drift-anchored (ours) & $\mathbf{0.791 \pm 0.013}$ & $0.547 \pm 0.023$ & 0.81 \\
margin     & $0.780 \pm 0.018$ & $0.706 \pm 0.015$ & 0.70 \\
qbc        & $0.779 \pm 0.019$ & $0.714 \pm 0.004$ & 0.70 \\
hybrid     & $0.779 \pm 0.019$ & $0.714 \pm 0.008$ & 0.71 \\
core-set   & $0.778 \pm 0.017$ & $0.687 \pm 0.012$ & 0.73 \\
badge      & $0.776 \pm 0.021$ & $0.708 \pm 0.012$ & 0.72 \\
random     & $0.762 \pm 0.016$ & $0.634 \pm 0.016$ & 0.87 \\
\midrule
\emph{ungated stream} & $0.766 \pm 0.011$ & $0.598 \pm 0.025$ & 0.84 \\
\bottomrule
\end{tabular}
\end{table}

The selection results then say something we did not expect. On PP, the informed
strategies are statistically indistinguishable --- they cluster between $0.776$
and $0.791$ --- so the choice of acquisition rule is \emph{not} where the
in-distribution gain comes from. Where it matters is the adversarial probe, and
there a clean trade-off appears: drift-anchored attains the best PP but the
\emph{worst} DP, while the more diffuse strategies give up about one PP point to
buy $14$--$17$ points of DP. Concentrating labels along the measured drift
sharpens a PP-specific boundary; spreading them out generalises to the
orthogonal adversarial shift. Two controls pin the cause down. Randomising
\emph{when} the loop fires, with the selector fixed, leaves PP essentially
unchanged ($0.798$ vs $0.800$), confirming that timing is not the lever; removing
the uncertainty pool and labelling the raw stream costs $2.5$ PP points,
confirming that the pool is. The trade-off itself is real and we revisit its
cause next: notably, the cleanest-label strategy (drift-anchored, oracle accuracy
$0.81$) is the OOD-weakest, a clue that the label \emph{noise} is doing useful
work.

\paragraph{Robustness to the one knob.}
\label{ssec:eval-alpha}
Drift-anchored carries a single weight, $\alpha$, on its drift term. Sweeping it
across five values shows the behaviour is a \emph{step}, not a tuned magnitude:
at $\alpha=0$ the selector is exactly margin sampling, turning the term on at any
positive weight reshuffles most of the selected set and produces the PP/DP
profile above, and increasing $\alpha$ further changes essentially nothing. The
specific value is immaterial --- it is the \emph{presence} of the drift term, not
its tuning, that defines the method (appendix).

\subsection{Choosing the Trigger: the Cost of Delay}
\label{ssec:eval-delay}

Timing is not our contribution, but the loop must fire on \emph{some} signal, and
the choice carries a measurable operational cost: between the onset of drift and
the first repair, the stale model keeps misclassifying in production. Replaying
the stream under whichever snapshot was live at each row turns that cost into a
number --- phishing missed before repair, and the time-integrated error
$\int(1{-}F_{1,w})$ --- and lets us compare triggers fairly
(Table~\ref{tab:al-delay}).

\begin{table}[H]
\centering
\caption{Cost of delay on $\mathrm{PP_{stream}}$ ($3{,}821$ rows, $1{,}865$
phishing; seed $2025$), replayed under the model snapshot live at each row.
``missed'' is phishing missed in production before repair; $\int(1{-}F_{1,w})$ is
the time-integrated operational error; ``probe'' counts oracle calls spent only
to \emph{measure} agreement, yielding no training label. We show the two
deployable label-free triggers, the output-gated reference, and two
representatives of the agreement-based (LAR) family; the full sweep is in the
appendix.}
\label{tab:al-delay}
\small
\begin{tabular}{@{}lccccc@{}}
\toprule
Trigger & 1st & \#fires & missed & $\int(1{-}F_{1,w})$ & probe \\
\midrule
\textsc{kl-only}            & \textbf{99}  & 30 & \textbf{828}  & \textbf{830.6}  & 0 \\
\textsc{periodic} (ours)    & 149 & 25 & 866  & 875.0  & 0 \\
\textsc{or} (output-gated)  & 429 & 2  & 1178 & 1217.4 & 0 \\
\midrule
LAR (best: random, $w{=}100$) & 299 & 17 & 1070 & 1103.3 & 380 \\
LAR (budget, $w{=}200$)       & 599 & 2  & 1144 & 1176.4 & 190 \\
\bottomrule
\end{tabular}
\end{table}

Two label-free triggers are deployable. \textsc{kl-only} watches the input
distribution and fires the moment it shifts: it reacts earliest (row $99$),
misses the fewest phishing, and incurs the lowest integrated error, at the cost
of firing often. \textsc{periodic} fires on a fixed schedule: it reacts a little
later but spends a \emph{predictable} budget, which is why we adopt it as the
canonical trigger --- a sensible default for a proof-of-concept, and within a
hair of \textsc{kl-only} on every metric ($0.74$ vs $0.75$ stream-wide weighted
$F_1$). The \textsc{or} trigger, which waits for the model's \emph{confidence} to
collapse, is the cautionary reference: on a calibrated detector that confidence
never collapses, so it fires twice in the entire stream and never leaves the
baseline.

We also asked whether the oracle itself could time the loop. The natural idea is
to periodically \emph{probe} a few samples with the LLM and fire when the model's
agreement with it drops --- a model-versus-oracle disagreement signal we call the
LLM-agreement rate (LAR), which we tried in several variants (random or
drift-guided probing, relative-drop or absolute-threshold rules, two window
sizes). It is intuitive but loses on every axis (Table~\ref{tab:al-delay}, lower
block): the best variant reacts $200$ rows later than \textsc{kl-only}, misses
over $240$ more phishing, and --- worse --- burns several hundred extra oracle
calls purely to measure agreement, calls that produce no label to train on. The
reason is structural: an agreement trigger needs the \emph{oracle} to witness the
drift, whereas the same shift is already visible, for free, in the model's own
prediction histogram. Watching the input is both the better signal and the
cheaper one, and it is what we deploy.

\subsection{Why an Imperfect Oracle Wins}
\label{sec:noise}

The trade-off above hinted that the oracle's mistakes are not pure cost. We now
show they are causally beneficial for robustness, with two experiments.

First, the ground-truth ceiling. Replacing the LLM with perfect labels,
everything else fixed, bounds what relabelling could ever buy. On PP it is a true
ceiling of $0.873$, of which our $\approx\!\$1$ oracle already captures two
thirds. But on the adversarial probe the perfect oracle is not a ceiling --- it
is a \emph{floor}: the ground-truth-trained model regresses to a macro $F_1$ of
$0.337$, \emph{below} the untrained baseline of $0.438$ and far below the LLM
run's $0.547$. Perfect labels sharpen the boundary so precisely onto PP's
phishing region that the $85\%$-benign probe gets swept up as phishing. If
perfect labels hurt OOD robustness and imperfect ones help it, the \emph{noise}
is doing the work.

Second, we isolate that claim. A naive LLM-versus-ground-truth comparison
confounds two things --- the labels differ and the selected samples can differ
--- so we hold the queried samples fixed and vary only label quality, flipping an
increasing fraction $\varepsilon$ of the oracle's mistakes back to truth and
refitting (Table~\ref{tab:cursor}). The trend is monotone and clean: as the
labels get \emph{better}, PP rises ($+9$ points) while every OOD metric falls
--- macro $F_1$, PR-AUC, phishing $F_1$ all collapse --- and the in-distribution
score stays flat, isolating the effect to the adversarial axis. At
$\varepsilon=1$ the cursor lands exactly on the independent ground-truth ceiling,
confirming the two methods agree.

\begin{table}[H]
\centering
\caption{Holding the queried samples fixed and correcting the oracle's errors
toward ground truth ($\varepsilon\!:\!0\to1$, margin selector, seed $2025$). As
labels improve, PP rises and the OOD probe falls monotonically; KP is unaffected.}
\label{tab:cursor}
\small
\begin{tabular}{@{}lccccc@{}}
\toprule
$\varepsilon$ & $\mathrm{PP}\,F_{1,w}$ & $\mathrm{PP}\,R_{\text{ph}}$ & $\mathrm{DP}\,F_{1,\mathrm{mac}}$ & $\mathrm{DP}\,\text{PR-AUC}$ & $\mathrm{KP}\,F_{1,w}$ \\
\midrule
0.00 & 0.775 & 0.589 & \textbf{0.683} & 0.753 & 0.973 \\
0.25 & 0.803 & 0.644 & 0.661 & 0.733 & 0.977 \\
0.50 & 0.834 & 0.705 & 0.613 & 0.709 & 0.978 \\
0.75 & 0.855 & 0.751 & 0.522 & 0.663 & 0.979 \\
1.00 & \textbf{0.866} & 0.812 & \textbf{0.350} & 0.568 & 0.978 \\
\bottomrule
\end{tabular}
\end{table}

The reading is direct: the oracle's $\approx\!29\%$ error rate acts as label
noise that keeps the boundary from over-committing to the PP-specific direction,
and so preserves robustness to an orthogonal adversarial shift; remove the noise
and the probe falls below its own untrained baseline. This reframes oracle
imperfection from a limitation into a mechanism --- in an adversarial deployment a
moderate oracle error rate is a feature, not a defect. It also raises a sharp
operational question for any deployment --- whether one should \emph{clean} the
oracle's labels before training on them --- which Section~\ref{sec:safety} takes
up, with a more nuanced answer than the noise result alone would suggest.
